\textbf{Bombeo y drenaje vigentes ENV-37/INT-37.} Se seleccionan dos bombas secundarias y sus accesorios en ENV-37; el circuito primario, fuente/serpentín, transferencia hidráulica, tolerancias de control y continuidad siguen pendientes. Los cálculos de ramal que siguen son sus antecedentes y no una segunda compra de bombas. INT-37 fija cuatro drenajes gravitacionales individuales y reserva de sifón negativo H350/J175/L557 mm con espacio vertical 650 mm; presión real, puerto/material, alarmas y cabida pendientes. El conjunto de recalentamiento mantiene sus funciones propias; selección secundaria no acredita frío húmedo ni recarga segura.

\section{Rama hidráulica y regulación del tratamiento de bancos: ENV-31}
\label{sec:ambiente-env31}
\textbf{Válvula y actuador seleccionados.} Se seleccionan cuatro válvulas Regin \texttt{ZTVB25-8}, una por cada ruta de tratamiento, con cuatro actuadores \texttt{RVAZ4-24A}. La válvula real es de dos vías, DN25, $K_{vs}=8$ m$^3$/h, PN16, carrera 5,5 mm y característica lineal; el diferencial máximo publicado para esta variante es 200 kPa. La ficha declara agua fría y medio de 1 a 110\,$^{\circ}$C, compatible con la propuesta de agua 5/10\,$^{\circ}$C. El cierre ocurre con vástago elevado; se respeta la flecha del cuerpo y no se instala el vástago hacia abajo. El fabricante identifica expresamente el RVAZ4-24A como actuador compatible. No se sustituye la característica lineal de ZTVB por la característica equiporcentual de ZTV ni se confunden las conexiones de ambas familias \parencite[pp.~1--2]{lote31-ztvb}.

El manual del actuador confirma 24 V AC/DC $\pm15\,\%$, control 0--10 V DC, potencia máxima 6 W, carrera 5,5 mm, fuerza 400 N y recorrido de 30 s. Se adopta alimentación DC desde la fuente protegida QUINT del camino correspondiente: dos conjuntos por riel ambiental de banco, cada riel respaldado desde A y B mediante ENV-40. Se conserva la reserva de cuatro actuadores, sin duplicarla al transferir su alimentación. Rojo recibe señal 0--10 V, negro positivo 24 V y azul referencia negativa, común con la salida analógica que lo controla. Se instalan en compartimento seco accesible de cada ruta, fuera de atmósfera de banco, con cuerpo motor hacia arriba o lateral dentro de 90 grados. Se usa ambiente de operación 0--50\,$^{\circ}$C y ausencia de condensación; IP44 no permite instalar la electrónica expuesta a humedad exterior ni aislarla cubriendo el motor \parencite[p.~1]{lote31-rvaz-manual}\parencite{lote31-rvaz-ficha}.

En régimen funcionan dos rutas y se reservan dos actuadores en movimiento: 12 W/0,5 A DC totales. Para transferencia, ejercicios y maniobra simultánea de los cuatro se fija máximo 24 W/1 A DC, 12 W/0,5 A por camino. La alimentación de las rutas en reserva permite cerrar/modular antes del cambio de ruta; no se descuenta su consumo por estar el ventilador parado. El balance de QUINT contempla el máximo de cuatro, sin sumar de nuevo 24 W como carga AC independiente. El calor de los actuadores se localiza en sus compartimentos; no se descuenta del recalentamiento del aire a menos que la implantación demuestre ese destino. No se modifica la reserva de impulsión ni los 18,1 kW máximos de calefactores/control.

\textbf{Dimensionamiento hidráulico propio.} Se conserva el caudal de aire de 900 m$^3$/h por banco. El siguiente cálculo hidráulico se refiere al punto de frío de 22,8604 kW y rocío 8,5 grados C; se retiene como contraste de ese estado, sin acreditarlo suficiente para la cota instrumental más exigente de ENV-34. La selección final debe recalcular caudal, pérdidas y autoridad al incorporar ese criterio, calor de trazado y serpentín real. Para agua sin glicol y salto 5 K, tomando propiedades de diseño $\rho=999,8$ kg/m$^3$ y $c_p=4,19$ kJ/(kg K), el caudal mínimo por ruta resulta
\[
q_w=\frac{22,8604}{999,8\cdot4,19\cdot5}\,3600=3,9291\;\mathrm{m^3/h}=1,0914\;\mathrm{l/s}.
\]
Dos rutas activas requieren 7,8581 m$^3$/h; el solape hidráulico de cuatro a pleno caudal exigiría 15,7163 m$^3$/h. La secuencia debe distinguir arranque de ventiladores de habilitación plena del agua de reserva: no se atribuye al circuito de dos rutas capacidad para cuatro por la sola presencia de dos bombas. Es caudal calculado a partir del mínimo de proceso; no es caudal OEM del serpentín pendiente ni garantiza su salida húmeda. Si la selección real exige otra potencia/salto, se recalculan válvula, autoridad, bomba y fuente antes de aceptar ese estado.

Con definición de $K_v$ para agua, la válvula completamente abierta tiene pérdida propia $\Delta p_v=100(q_w/8)^2=24,1212$ kPa. No se identifica esa pérdida con el diferencial máximo de catálogo de 200 kPa. Se adopta como criterio de regulación de obra una autoridad mínima 0,30 y se limita el resto de pérdidas del ramal a 55 kPa: $a=24,1212/(24,1212+55)=0,3049$, diferencial de ramal 79,1212 kPa. Ese diferencial de diseño queda bajo el límite de la válvula; el máximo de bomba a caudal nulo y en ramas cerradas también debe permanecer bajo 200 kPa mediante regulación/limitación independiente. No se afirma que exista ya una bomba que satisfaga ambas condiciones.

Como contraste de trazado se conserva una longitud equivalente de tubería recta máxima de 40 m por rama, diámetro interior de proyecto mínimo de 40 mm, rugosidad de cálculo 0,0015 mm y suma de pérdidas singulares $\sum K=12$. Son cotas de obra de ida y retorno, no medidas del edificio existente ni catálogo de una tubería seleccionada. A 3,9291 m$^3$/h, $v=0,8685$ m/s y, con viscosidad de diseño 1,385 mPa s, $Re=25.078$. Usando aproximación de Haaland se obtiene $f=0,024409$ y Darcy--Weisbach da
\[
\Delta p_{tub}=\left(f\frac{40}{0,040}+12\right)\frac{999,8(0,8685)^2}{2}=13,7290\;\mathrm{kPa}.
\]
Quedan 41,2710 kPa de la cota de 55 para serpentín, filtro hidráulico, válvula de equilibrado y accesorios no incluidos en los doce $K$. El DN25 de la válvula no reduce todo el circuito a DN25: sus conexiones y transiciones se concentran en la estación de válvula y su pérdida está separada. El presupuesto se acepta sólo si las curvas reales de esos accesorios y serpentín respetan la suma. La bomba se seleccionará contra el ramal más desfavorable más distribución común, intercambiador/fuente y accesorios generales; 79,1212 kPa no es altura de bomba completa y no se usa para inventar potencia/arranque del motor.

\textbf{Regulación y pérdida de energía.} La válvula es elemento de regulación de frío y no dispositivo de seguridad contra H$_2$, incendio o fuga de agua. La ficha declara ausencia de fuga con cierre, pero el actuador no acredita retorno por resorte ni posición segura al perder DC. Sus 30 s son recorrido nominal de regulación, no tiempo de corte de recarga ni garantía de aislamiento. Al perder señal o alimentación se retiran calor y permiso de recarga por la cadena independiente; no se espera a que la válvula alcance cierre. La retirada de agua por fuga requiere elemento final y supervisión independientes aún pendientes. El enfriamiento se modula por rocío del aire tratado, el recalentamiento por temperatura final y la velocidad EC por caudal; son tres lazos distintos. ENV-34 selecciona los nodos ordinarios WAGO, AI 750-455, AO 750-559 y cuatro DMT132; asigna una salida individual a cada actuador y desarrolla saturación, antioscilación y mapa de señales. Falta verificar la impedancia real de cada entrada RVAZ4-24A frente a la carga mínima de 5 k$\Omega$ del AO, el sentido eléctrico de apertura/cierre y las ganancias PI sobre el conjunto húmedo real. El sensor de rocío y sus tiempos típicos no habilitan recarga; auxiliares e interfaces se concilian con la cadena independiente de seguridad. La indicación manual 0--100\,\% del actuador no constituye contacto eléctrico de posición ni prueba de caudal de agua.

\textbf{Frontera del serpentín e implantación.} El catálogo primario CWK/CFK verifica variantes de tres filas, dimensiones, bandeja y necesidad de selección con caudal, temperatura, humedad, agua y glicol; no entrega capacidad húmeda del punto 35\,$^{\circ}$C/80\,\% HR, 900 m$^3$/h y agua 5/10\,$^{\circ}$C. Por ello no se adopta CWK~315-3-2.5 ni CFK~315-3-2.5 como sustituto de la batería a medida por su diámetro. La CFK~315 mide 639 por 553 mm y longitud 462 mm, con aislamiento de 50 mm: esas dimensiones reales permiten cribar instalación, pero no prueban capacidad o factor de bypass. El manual exige conducto horizontal, inclinación 10$\pm1$ grados hacia drenaje para esta familia y evitar cercanía a descarga de ventilador/codo. No se transfieren esas instrucciones a un serpentín de otro fabricante. La fuente Carrier corregida, batería húmeda, pérdidas de aire, rango mínimo, condensado y bomba permanecen pendientes de conjunto concreto \parencite[pp.~2--5]{lote31-cwk-catalogo}\parencite[p.~3]{lote31-cwk-manual}.

Se mantienen las cuatro reservas de UTA de 3 por 2 m, plenums ENV-25 y estaciones adicionales ENV-28 sobre bases de 2,5 por 1 m. Las válvulas se disponen en el retorno de cada ruta, con unión desmontable, purga, vaciado y acceso de obra; no se atribuye cabida completa por el tamaño del actuador. El plano debe mostrar conexiones sin esfuerzo sobre serpentín, drenaje separado, aislamiento continuo y acceso a bandeja, filtro y válvula, dentro de las reservas vigentes antes de cerrar RT-06.03. Este desarrollo concreta modelo, compatibilidad, mando, cantidad y demanda hidráulica de regulación. RT-06.03, RT-06.09 y RT-06.13 conservan estado pendiente hasta integrar selección húmeda y fuente, bombas, pérdidas completas, implantación y seguridad.

El detalle de aislamiento distingue válvula/tubería fría expuesta al exterior de los interiores tratados: INT-34 fija al menos 100 mm de reserva de obra exterior y resolución de puentes, sin cubrir el motor del actuador ni prohibir su ventilación. El espesor de 25 mm de ENV-25 sólo corresponde a entorno interior tratado a 24 grados C. Las dimensiones de cuerpo, conexiones y acceso se concilian con el revestimiento exterior real antes de declarar encaje.
